\appendix
\section{Blocs à copier-coller}\label{sec:blocs}
Tous les blocs ci-dessous sont générés depuis les fichiers de \code{data/} par \code{tools/make\_blocks.py} : ce qui est collé est identique au fichier. Les tabulations se perdent en copiant depuis un PDF, d'où le choix « une colonne par bloc » ou « fichier à coller en une fois ».
\subsection{Sources de référence : une colonne par bloc (\code{A-1} à \code{A-12})}\label{sec:blocsA}
Sélectionner un bloc, copier, cliquer sur la \textbf{première cellule} de la colonne correspondante dans \ui{Radiating Sources}, puis \code{Ctrl+V}. Les sept lignes sont \code{TS01}\ldots\code{TS07}.
\noindent
\begin{minipage}[t]{0.325\textwidth}\small\textbf{A-1} Radiation source name\\[-2pt]
\begin{Verbatim}[frame=single,framesep=1mm,fontsize=\footnotesize,rulecolor=\color{accent}]
TS01_900_man
TS02_2100_3tx
TS03_23GHz_PtP
TS04_18GHz_PtP
TS05_LTE_att
TS06_ISO_1W
TS07_ISO_1kW
\end{Verbatim}
\end{minipage}\hfill
\begin{minipage}[t]{0.325\textwidth}\small\textbf{A-2} Sector ID\\[-2pt]
\begin{Verbatim}[frame=single,framesep=1mm,fontsize=\footnotesize,rulecolor=\color{accent}]
A
B
C
C
A
-
-
\end{Verbatim}
\end{minipage}\hfill
\begin{minipage}[t]{0.325\textwidth}\small\textbf{A-3} Type of service\\[-2pt]
\begin{Verbatim}[frame=single,framesep=1mm,fontsize=\footnotesize,rulecolor=\color{accent}]
GSM
UMTS
PtP
PtP
LTE
Other
Other
\end{Verbatim}
\end{minipage}\hfill
\par\medskip
\noindent
\begin{minipage}[t]{0.325\textwidth}\small\textbf{A-4} Frequency, MHz\\[-2pt]
\begin{Verbatim}[frame=single,framesep=1mm,fontsize=\footnotesize,rulecolor=\color{accent}]
900
2100
23000
18000
2600
900
900
\end{Verbatim}
\end{minipage}\hfill
\begin{minipage}[t]{0.325\textwidth}\small\textbf{A-5} Transmitter power, W\\[-2pt]
\begin{Verbatim}[frame=single,framesep=1mm,fontsize=\footnotesize,rulecolor=\color{accent}]
12.05
10.42
0.1995
0.2512
20
1
1000
\end{Verbatim}
\end{minipage}\hfill
\begin{minipage}[t]{0.325\textwidth}\small\textbf{A-6} Number of TX\\[-2pt]
\begin{Verbatim}[frame=single,framesep=1mm,fontsize=\footnotesize,rulecolor=\color{accent}]
1
3
1
1
1
1
1
\end{Verbatim}
\end{minipage}\hfill
\par\medskip
\noindent
\begin{minipage}[t]{0.325\textwidth}\small\textbf{A-7} Total atten., dB\\[-2pt]
\begin{Verbatim}[frame=single,framesep=1mm,fontsize=\footnotesize,rulecolor=\color{accent}]
0
0
0
0
0.37
0
0
\end{Verbatim}
\end{minipage}\hfill
\begin{minipage}[t]{0.325\textwidth}\small\textbf{A-8} Antenna model\\[-2pt]
\begin{Verbatim}[frame=single,framesep=1mm,fontsize=\footnotesize,rulecolor=\color{accent}]
742265V02
742265V02
23G 0.6
18G 0.6
AMB4519
isotropic
isotropic
\end{Verbatim}
\end{minipage}\hfill
\begin{minipage}[t]{0.325\textwidth}\small\textbf{A-9} Antenna azimuth, deg\\[-2pt]
\begin{Verbatim}[frame=single,framesep=1mm,fontsize=\footnotesize,rulecolor=\color{accent}]
60
130
125
220
180
0
0
\end{Verbatim}
\end{minipage}\hfill
\par\medskip
\noindent
\begin{minipage}[t]{0.325\textwidth}\small\textbf{A-10} Antenna gain, dBi\\[-2pt]
\begin{Verbatim}[frame=single,framesep=1mm,fontsize=\footnotesize,rulecolor=\color{accent}]
16
18.3
40
39.2
18
0
0
\end{Verbatim}
\end{minipage}\hfill
\begin{minipage}[t]{0.325\textwidth}\small\textbf{A-11} Antenna height, m\\[-2pt]
\begin{Verbatim}[frame=single,framesep=1mm,fontsize=\footnotesize,rulecolor=\color{accent}]
29
29
27.2
27.2
40
10
10
\end{Verbatim}
\end{minipage}\hfill
\begin{minipage}[t]{0.325\textwidth}\small\textbf{A-12} Antenna size, m\\[-2pt]
\begin{Verbatim}[frame=single,framesep=1mm,fontsize=\footnotesize,rulecolor=\color{accent}]
1.9
1.9
0.6
0.6
1.5
0.1
0.1
\end{Verbatim}
\end{minipage}\hfill
\par\medskip
\begin{note}Pour \code{TS07} (1~kW isotrope), l'oracle des contours circulaires (\code{HCT-04}) suppose une \textbf{hauteur d'antenne de 10~m} et le diagramme \code{ok\_isotropic.msi}. Pour \code{BLD-11}, passez-la à 20~m.\end{note}
\subsection{Limites, origine, formats}\label{sec:blocsC}
\paragraph*{Bloc \code{C-1} --- limites plates (9 valeurs, de gauche à droite dans la ligne 1 du tableau \ui{EMF Limits})}
\begin{center}\small\begin{tabular}{@{}ccccccccc@{}}\toprule \textbf{Freq.\ to F} & \textbf{E1} & \textbf{A1} & \textbf{E2} & \textbf{A2} & \textbf{E3} & \textbf{A3} & \textbf{E4} & \textbf{A4}\\\midrule 1000000 & 6 & 0 & 20 & 0 & 4.5 & 0 & 10 & 0\\\bottomrule\end{tabular}\end{center}
Fichier à coller en une fois : \code{data/limits/limits\_flat.tsv}.
\paragraph*{Bloc \code{C-2} --- chaînes à coller pour l'origine (une ligne à la fois, \code{ORI-02})}\vspace{-2pt}
\begin{Verbatim}[frame=single,framesep=1.5mm,fontsize=\footnotesize,rulecolor=\color{accent}]
45.485337, 9.198818
45.485337 9.198818
45,485337; 9,198818
N45.485337 E9.198818
45°29'07.2"N 9°11'55.7"E
45°29.1202'N 9°11.9291'E
   45.485337 ,   9.198818   
9.198818, 45.485337
91.0, 9.198818
45.485337, 181.0
abc, def
45.485337
\end{Verbatim}
\paragraph*{Bloc \code{C-3} --- table ICNIRP du manuel (p.21), 4 lignes}
\begin{center}\small\begin{tabular}{@{}rrrrr@{}}\toprule \textbf{Freq.\ to F} & \textbf{E1} & \textbf{A1} & \textbf{E2} & \textbf{A2}\\\midrule 30 & 824 & $-1$ & 1842 & $-1$\\ 300 & 27.5 & 0 & 61.4 & 0\\ 1500 & 1.585 & 0.5 & 3.545 & 0.5\\ 1000000 & 61.4 & 0 & 137.3 & 0\\\bottomrule\end{tabular}\end{center}
Fichier à coller : \code{data/limits/limits\_icnirp.tsv}.\\ Valeurs attendues de $E_L$ aux fréquences de test : \code{data/limits/limit\_check\_icnirp.tsv}.
\paragraph*{Bloc \code{C-4} --- aller-retour des formats de coordonnées (\code{ORI-03})}
\begin{center}\scriptsize\renewcommand{\arraystretch}{1.1}\begin{tabular}{@{}rrp{4.6cm}p{5.2cm}@{}}\toprule \textbf{lat} & \textbf{lon} & \textbf{DMS attendu} & \textbf{DM attendu}\\\midrule
45.485337 & 9.198818 & 45°29'07.21'' N 9°11'55.74'' E & 45°29.1202' N 9°11.9291' E\\
0 & 0 & 0°00'00.00'' N 0°00'00.00'' E & 0°00.0000' N 0°00.0000' E\\
-33.8688 & 151.2093 & 33°52'07.68'' S 151°12'33.48'' E & 33°52.1280' S 151°12.5580' E\\
5.36 & -4.0083 & 5°21'36.00'' N 4°00'29.88'' W & 5°21.6000' N 4°00.4980' W\\
90 & 180 & 90°00'00.00'' N 180°00'00.00'' E & 90°00.0000' N 180°00.0000' E\\
-90 & -180 & 90°00'00.00'' S 180°00'00.00'' W & 90°00.0000' S 180°00.0000' W\\
10.999999 & 10.999999 & 11°00'00.00'' N 11°00'00.00'' E & 10°59.9999' N 10°59.9999' E\\
44.345678 & -134.567893 & 44°20'44.44'' N 134°34'04.41'' W & 44°20.7407' N 134°34.0736' W\\
\bottomrule\end{tabular}\end{center}
\subsection{Textes et nombres pour la robustesse}\label{sec:blocsF}
\paragraph*{Bloc \code{F-1} --- chaînes de texte (une par ligne ; les lignes vides, très longues et multi-scripts sont dans \code{fuzz/strings\_text.txt})}\vspace{-2pt}
\begin{Verbatim}[frame=single,framesep=1.5mm,fontsize=\footnotesize,rulecolor=\color{accent}]
   leading and trailing   
quote ' and " and `
semi;colon;list
back\slash C:\temp\x
../../etc/passwd
<script>alert(1)</script>
&amp; <b>bold</b> &lt;
%s %d %n {0} {1}
{SiteName}
{RadiatingSources_1}
{{nested}}
'; DROP TABLE sources;--
=cmd|' /C calc'!A0
CON
NUL.txt
0÷360
\end{Verbatim}
Chaînes longues à fabriquer au besoin : \code{python3 -c "print('A'*5000)"} (5000 caractères). Pour l'Unicode multi-scripts, utilisez le fichier : les polices du PDF ne les contiennent pas.
\paragraph*{Bloc \code{F-2} --- nombres écrits de façon piégeuse (\code{ROB-02})}\vspace{-2pt}
\begin{Verbatim}[frame=single,framesep=1.5mm,fontsize=\footnotesize,rulecolor=\color{accent}]
0
-0
-1
+5
5.
  .5
1e3
1E-3
1,5
1.5
1 500
1'500
1.2.3
--5
1..2
0x10
12abc
abc
NaN
Infinity
-Infinity
1e400
1e-400
2147483647
2147483648
4294967296
9223372036854775808
0.1+0.2
3,14159
3.14159265358979323846264338327950288
\end{Verbatim}
\paragraph*{Bloc \code{F-3} --- bornes à essayer par champ (\code{fuzz/boundaries.tsv})}
\begin{center}\scriptsize\renewcommand{\arraystretch}{1.1}\begin{tabular}{@{}lp{12.2cm}@{}}\toprule \textbf{Champ} & \textbf{Valeurs} \\\midrule
frequency\_MHz & 0 | 0.001 | 0.099999 | 0.1 | 300000 | 300000.001 | 1000000 | 1000000.001 | -900\\
power\_W & 0 | 1e-12 | 1e6 | 1e9 | -1\\
number\_of\_TX & 0 | 1 | 1.5 | 1000 | -1\\
total\_atten\_dB & -10 | 0 | 100 | 1000\\
gain\_dBi & -30 | 0 | 60 | 200\\
height\_m & 0 | -5 | 0.01 | 1000 | 100000\\
azimuth\_deg & 0 | 360 | 360.0001 | -1 | 720 | x (non numeric -> '0÷360' per manual p.16)\\
mech\_tilt\_deg & -90 | -91 | 0 | 90 | 91\\
reflection\_coef & 0.99 | 1 | 1.3 | 1.6 | 1.61 | 0 | -1 | 1,3\\
matrix\_size & 0 | 49 | 50 | 800 | 1000 | 1001 | 100000 | abc\\
grid\_dimension\_m & 0 | -100 | 0.5 | 100 | 100000 | 1e9\\
building\_height\_m & 0 | -5 | 3 | 12.5 | 10000\\
\bottomrule\end{tabular}\end{center}
\paragraph*{Bloc \code{L-1} --- étiquettes de tour (\code{LBL-03}, \code{LBL-04})}\vspace{-2pt}
\begin{Verbatim}[frame=single,framesep=1.5mm,fontsize=\footnotesize,rulecolor=\color{accent}]
Tower/15/15/0.5/1.2
Tower/15/0/0.5/0.5
Tower/15/15/0.5
tower/15/15/0.5/1.2
Tower/-5/15/0.5/1.2
Tower/0/0/0/0
Tower/15/15/0,5/1,2
Tower / 15 / 15 / 0.5 / 1.2
Tower/abc/15/0.5/1.2
Tower/1e3/1e3/5/5
Tower////
\end{Verbatim}
\paragraph*{Bloc \code{R-1} --- alias personnalisés (\code{RPT-03}) ; fichier \code{aliases/custom\_aliases.tsv}}
\begin{center}\scriptsize\renewcommand{\arraystretch}{1.1}\begin{tabular}{@{}lp{11.3cm}@{}}\toprule \textbf{Alias name} & \textbf{Alias text}\\\midrule
\{SiteName\} & TEST SITE é/ü [emoji]\\
\{ProjectName\} & Project <A\&B> "quoted"\\
\{SiteAddress\} & Long address Long address Long\ldots{} (80 répétitions de « Long address », voir fichier)\\
\{Customer\} & ACME;Corp\\
\{Designer\} & \emph{(vide)}\\
\{Date\} & 2026-10-20\\
\{RadiatingSources\_1\} & collision with standard alias\\
\{Unused\} & never used in template\\
\{SiteName\} & duplicate name\\
\{ with space \} & space in name\\
\bottomrule\end{tabular}\end{center}
\paragraph*{Bloc \code{S-1} --- chaînes pour le champ d'activation (\code{INS-05}) --- ne servent qu'à vérifier le refus}\vspace{-2pt}
\begin{Verbatim}[frame=single,framesep=1.5mm,fontsize=\footnotesize,rulecolor=\color{accent}]

     
12345
'";--
AAAA-AAAA-AAAA-AAAA
0000-0000-0000-0000
\end{Verbatim}
\paragraph*{Bloc \code{W-1} --- limiter les c\oe urs de calcul (\code{ROB-06}), invite de commandes Windows ; adaptez le chemin}\vspace{-2pt}
\begin{Verbatim}[frame=single,framesep=1.5mm,fontsize=\footnotesize,rulecolor=\color{accent}]
start "" /affinity 1  "C:\Chemin\vers\EMFPlanner.exe"
start "" /affinity 3  "C:\Chemin\vers\EMFPlanner.exe"
start "" /affinity F  "C:\Chemin\vers\EMFPlanner.exe"
start "" /affinity FF "C:\Chemin\vers\EMFPlanner.exe"
\end{Verbatim}
\paragraph*{Bloc \code{P-1} --- commandes de vérification (Python 3, depuis \code{emf\_test\_kit/})}\vspace{-2pt}
\begin{Verbatim}[frame=single,framesep=1.5mm,fontsize=\footnotesize,rulecolor=\color{accent}]
python3 tools/gen_all.py
python3 tools/msi_check.py sortie.msi 7000 31 0.6
python3 tools/msi_check.py sortie.msi 40000 50 2.0
python3 tools/env_check.py enveloppe.msi data/msi/envelope/*.msi
python3 tools/tile_server.py
\end{Verbatim}


\section{Lecture des cas limites du KML}\label{sec:kmledge}
\code{buildings\_edge.kml} contient 19 éléments. Pour chacun, relevez : importé ? hauteur lue ? libellé lu ? position correcte ?
\begin{center}\scriptsize\renewcommand{\arraystretch}{1.12}
\begin{tabular}{@{}rp{3.6cm}p{5.1cm}p{5.2cm}@{}}\toprule
\textbf{\#} & \textbf{Nom dans le KML} & \textbf{Particularité} & \textbf{À regarder}\\\midrule
1 & \code{8} & pas de point-virgule & hauteur 8~? libellé vide ? ou refus ?\\
2 & \code{;NoHeight} & hauteur absente & hauteur par défaut ou 0 ?\\
3 & \code{abc;NotNumber} & hauteur non numérique & refus, 0 ou valeur silencieuse ?\\
4 & \code{0;Zero} & hauteur nulle & bâtiment plat, peut-il bloquer ?\\
5 & \code{-5;Negative} & hauteur négative & refus attendu\\
6 & \code{10000;Huge} & 10~km de haut & plafonné ? 3D déformée ?\\
7 & \code{12,5;CommaDecimal} & virgule décimale & 12,5 ou 125 ou 12 ?\\
8 & \code{9;Église é ü [CJK] [coche]} & Unicode & libellé conservé (le PDF n'affiche pas les idéogrammes)\\
9 & \code{9;a;b;c} & plusieurs points-virgules & libellé \code{a;b;c} ou \code{a} ?\\
10 & \code{9;Hole} & polygone avec trou intérieur & trou ignoré ou respecté ?\\
11 & \code{9;Collinear} & 3 points alignés & rejet attendu\\
12 & \code{9;Bowtie} & polygone croisé (nœud papillon) & rejet ou rendu cohérent\\
13 & \code{10;B01} & doublon de \code{buildings\_ok.kml} & doublon créé ?\\
14 & \code{5;Tiny} & 0,5~m de côté & importé sans plantage\\
15 & \code{20;Far10km} & à 14~km de l'origine & import OK ; 3D impossible tant qu'il reste (\code{HCT-14})\\
16 & \code{9;NoAltitude} & coordonnées sans altitude (\code{lon,lat}) & importé normalement\\
17 & \code{9;Unclosed} & anneau non fermé & fermé automatiquement ou rejeté\\
18 & \code{7;PointNotPolygon} & \code{Point} & ignoré, aucun bâtiment\\
19 & \code{7;LineNotPolygon} & \code{LineString} & ignoré, aucun bâtiment\\\bottomrule
\end{tabular}\end{center}

\section{Formules F.699-8 utilisées comme oracle}\label{sec:f699}
Avec $\lambda=299{,}792458/f_{\mathrm{MHz}}$ (m), $D/\lambda$ rapport diamètre/longueur d'onde, $\varphi$ angle hors axe (degrés), $G_{\max}$ gain maximal (dBi) :\\[2pt]
\textbf{Cas $D/\lambda\le100$} : $G_1=2+15\log_{10}(D/\lambda)$ ; $\varphi_m=\dfrac{20\lambda}{D}\sqrt{G_{\max}-G_1}$.
\[G(\varphi)=\begin{cases}G_{\max}-2{,}5\cdot10^{-3}\left(\frac{D\varphi}{\lambda}\right)^2&0\le\varphi<\varphi_m\\ G_1&\varphi_m\le\varphi<100\lambda/D\\ 52-10\log_{10}\frac D\lambda-25\log_{10}\varphi&100\lambda/D\le\varphi<48^\circ\\ 10-10\log_{10}\frac D\lambda&48^\circ\le\varphi\le180^\circ\end{cases}\]
\textbf{Cas $D/\lambda>100$} : $G_1=-1+15\log_{10}(D/\lambda)$ ; $\varphi_r=15{,}85\,(D/\lambda)^{-0{,}6}$.
\[G(\varphi)=\begin{cases}G_{\max}-2{,}5\cdot10^{-3}\left(\frac{D\varphi}{\lambda}\right)^2&0\le\varphi<\varphi_m\\ G_1&\varphi_m\le\varphi<\varphi_r\\ 29-25\log_{10}\varphi&\varphi_r\le\varphi<48^\circ\\ -13&48^\circ\le\varphi\le180^\circ\end{cases}\]
Atténuation relative écrite dans le MSI : $G_{\max}-G(\varphi)$. Les valeurs de référence sont dans \code{data/msi/f699\_expected\_*.tsv}. Continuité contrôlée aux frontières de zone (écart $<0{,}05$~dB hors $\varphi_m$ où le plateau commence par $G_{\max}-(G_{\max}-G_1)=G_1$ exactement).
\begin{alerte}Ces formules sont reconstituées de mémoire : relisez-les contre le texte de la Rec.\ UIT-R F.699-8 avant de conclure à un défaut.\end{alerte}

\section{Feuille de résultats et fiche de défaut}\label{sec:fiche}
\subsection*{Feuille de résultats (à recopier dans un tableur)}
\begin{center}\small\begin{tabular}{@{}|p{1.5cm}|p{1.9cm}|p{1.5cm}|p{2.1cm}|p{2.2cm}|p{2.1cm}|p{2.8cm}|@{}}\hline
\textbf{ID} & \textbf{Date} & \textbf{Version} & \textbf{Valeur lue} & \textbf{Valeur attendue} & \textbf{OK/KO/N/A} & \textbf{Preuve (fichier)}\\\hline\ \\[3mm]\hline\ \\[3mm]\hline\end{tabular}\end{center}
\subsection*{Fiche de défaut}
\begin{Verbatim}[frame=single,framesep=2mm,fontsize=\footnotesize,rulecolor=\color{accent}]
Titre            : (une phrase : quoi, où)
ID du test       : ex. SRC-13            Gravité : bloquant / majeur / mineur / cosmétique
Version/OS/lang  : EMF Planner 7.1 (build ...), Windows 11, région FR/EN, Office ...
Étapes           : 1) ...  2) ...  3) ...
Attendu (oracle) : valeur + [M]/[F]/[R]/[X]/[O] + source (page du manuel ou fichier du kit)
Observé          : valeur + capture
Reproductible    : toujours / parfois (n sur 5) / une fois
Fichiers joints  : .xemf, capture, données collées, message d'erreur exact
\end{Verbatim}

\section{Glossaire anglais $\leftrightarrow$ français}
Les libellés du logiciel sont en anglais ; ce petit tableau aide à lire le manuel et les messages d'erreur.
\begin{center}\small\renewcommand{\arraystretch}{1.1}
\begin{tabular}{@{}ll@{\qquad}ll@{}}\toprule
\textbf{English} & \textbf{Français} & \textbf{English} & \textbf{Français}\\\midrule
radiating source & source rayonnante & feeder & câble d'alimentation\\
jumper & jarretière & attenuation & affaiblissement\\
gain & gain & azimuth & azimut\\
mechanical / electrical tilt & inclinaison mécanique / électrique & beamwidth & largeur de faisceau\\
antenna pattern & diagramme d'antenne & envelope & enveloppe\\
exposure limit & limite d'exposition & field strength & intensité de champ\\
safety volume & volume de sécurité & control point & point de contrôle\\
horizontal / vertical cut & coupe horizontale / verticale & calculation grid & grille de calcul\\
base map & fond de carte & tile & tuile\\
shielding & blindage & contour line & courbe de niveau\\
trusted location & emplacement approuvé & alias & alias (signet du modèle Word)\\
reflection coefficient & coefficient de réflexion & duty cycle & rapport cyclique\\\bottomrule
\end{tabular}\end{center}
